\section{Conclusion}

We set out to test whether adding mypy to LLM repair loops improves pass@1.
It does --- by nearly 9 percentage points, consistently across three seeds. We
also set out to understand why. The answer surprised us: the improvement is
not where we expected it.

The type-specificity mechanism --- mypy identifies a type error, tells the LLM
exactly what went wrong, the LLM corrects it specifically --- turns out not to
be the active channel at $k$=5 saturation. Both conditions achieve identical
90.9\% repair rates on type-error problems. The +8.94pp overall improvement
arises elsewhere, most plausibly from general diagnostic context enrichment:
mypy output adds structured text to repair prompts that improves GPT-4o-mini's
repair quality across all problem types.

Our contributions are: (1) first controlled evidence that execution+mypy
outperforms execution-only repair on HumanEval+ (+8.94pp, three seeds
consistent); (2) pre-registered type-specificity refutation
(differential=0.000, $n$=22); (3) benchmark-level mypy signal characterization
(70\% vs 0\%); and (4) Z3 spec generation pipeline (91.1\% validity). Together,
these results establish both the \emph{what} (mypy helps substantially) and the
\emph{what not} (not via type-targeted correction), opening a research agenda
about what makes feedback signals distinctive from execution output.

\textbf{Future directions} include: (a) Condition B-Null ablation to distinguish
mypy content from prompt length effects (highest priority); (b) per-round
type-specificity analysis at $k$=1..2 before ceiling saturation; (c)
multi-model replication (GPT-4o, open-source code models); and (d) Z3 CE
feedback on harder benchmarks (LiveCodeBench) with lower base pass@1 and more
headroom.

We began by asking whether adding mypy was worth the one-subprocess-call
investment. The answer is yes. The research question that emerged is richer:
\emph{why} does mypy help, and where does its benefit concentrate?
Understanding what makes a feedback signal distinctive from execution output ---
whether it is error specificity, structured format, or prompt enrichment --- is
the key to principled design of LLM repair loop architectures.
